\todo{Lecture 17}
\chapter{Développements Limites}
\begin{definition}
Soit $f:D\to \mathbb{R}$ une fonction $n$ fois derivable en $x_{0}$. Le polynome de Taylor de $f$ en $x_{0}$ a l'ordre $n$ est
$$
p_{n}(x)=\sum_{k=0}^{n}  \frac{f^{(k)}(x_{0})}{k!}(x-x_{0})^{k},
$$
avec les conventions $f^{(0)}=f$, $0! =1$ et $(x-x_{0})^{0}=1$.
\end{definition}
\begin{remark}
Par construction, $p_{n}$ est l'unique polynome de degre au plus $n$ tel que $f^{(n)}(x_{0})=p_{n}^{(n)}(x_{0})$.
\end{remark}
\begin{theorem}[Formule de Taylor-Lagrange]
Soient $n\in \mathbb{N}$, $I$ un intervalle ouvert et $f:I\to \mathbb{R}$ une fonction $n+1$ fois derivable sur $I$. Fixons $x_{0}\in I$. Pour tout $x \in I\setminus \{ x_{0} \}$, il existe $c_{x}$ strictement entre $x_{0}$ et $x$ tel que
$$
f(x)=p_{n}(x)+\frac{f^{(n+1)}(c_{x})}{(n+1)!}(x-x_{0})^{n+1}.
$$
\end{theorem}
\begin{remark}
La difference $r_{n}(x):=f(x)-p_{n}(x)$ s'appelle le reste a l'ordre $n$.
\end{remark}
\begin{proof}
Fixons $x\neq x_{0}$. Considerons
$$
s(t)=\sum_{k=0}^{n}  \frac{f^{(k)}(t)}{k!}(x-t)^{k}.
$$
On a donc $s(x)=f(x)$ et $s(x_{0})=p_{n}(x)$. Definissons
$$
h(t)=s(t)+\lambda(x-t)^{n+1}
$$
avec $\lambda \in \mathbb{R}$ a choisir tel que $h(x)=h(x_{0})$. Donc
$$
\lambda= \frac{f(x)-p_{n}(x)}{(x-x_{0})^{n+1}}.
$$
On peut appliquer Rolle a $h$ en $x$ et $x_{0}$ car $h$ est derivable sur $I$ et $h(x)=h(x_{0})$. Donc, il existe $c_{x}$ strictement entre $x$ et $x_{0}$ tel que $0=h'(c_{x})$. On a
\begin{align*}
s'(t) & =f'(t)+\sum_{k=1}^{n} \left[ \frac{f^{(k+1)}(t)}{k!}(x-t)^{k}-\frac{f^{(k)}(t)}{k!}k(x-t)^{k} \right] \\
 & =f'(t)+\sum_{k=1}^{n} \frac{f^{(k+1)}(t)}{k!}(x-t)^{k}-\sum_{k=1}^{n} \frac{f^{k}(t)}{(k-1)!}(x-t)^{k-1} \\
 &  = f'(t)+\sum_{k=1}^{n}  \frac{f^{(k+1)}(t)}{k!}(x-t)^{k}-\sum_{k=0}^{n-1} \frac{f^{(k+1)}(t)}{k!}(x-t)^{k} \\
 &   =f'(t)+\frac{f^{n+1}(t)}{n!}(x-t)^{n}-f'(t) \\
 &   = \frac{f^{(n+1)}(t)}{n!}(x-t)^{n}.
\end{align*}
Donc
\begin{align*}
0 & =h'(c_{x}) \\
 & =s'(c_{x})-\lambda(n+1)(x-c_{x})^{n} \\
 &  = \frac{f^{(n+1)}(c_{x})}{n!}(x-c_{x})^{n}-\lambda(n+1)(x-c_{x})^{n}.
\end{align*}`'
Comme $c_{x}\neq x$, on a
$$
\frac{f(x)-p_{n}(x)}{(x-x_{0})^{n+1}}=\lambda=\frac{f^{(n+1)}(c_{x})}{(n+1)!}.
$$
En reorganisant, on trouve
$$
f(x)-p_{n}(x)=\frac{f^{(n+1)}(c_{x})}{(n+1)!}(x-x_{0})^{n+1}.
$$
\end{proof}
\begin{corollary}
Dans le contexte de Taylor-Lagrange, s'il existe $M\ge 0$ tel que
$$
	|f^{(n+1)}(t)|\le M,\quad \forall t\in I,
$$
alors, pour tout $x \in I$, 
$$
|f(x)-p_{n}(x)|\le \frac{M}{(n+1)!}|x-x_{0}|^{n+1}.
$$
\end{corollary}
\begin{theorem}[Formule de Taylor-Young]
Soient $n\in \mathbb{N}^{*}$, $I$ un intervalle ouvert et $f:I\to \mathbb{R}$ une fonction $n-1$ fois derivable sur $I$. Fixons $x_{0}\in I$ et supposons que $f^{(n-1)}$ est derivable en $x_{0}$.

Alors, il existe $\varepsilon:I\to \mathbb{R}$ telle que, pour tout $x \in I$,
$$
f(x)=p_{n}(x)+\varepsilon(x)(x-x_{0})^{n},
$$
avec $\varepsilon(x_{0})=0$ et
$$
\lim_{ x \to x_{0} } \varepsilon(x)=0.
$$
\end{theorem}
\begin{proof}[Preuve d'une version simplifiee de Taylor-Young]
Supposons que $f$ est $n$ fois derivable sur $I$ et que $f^{(n)}$ est continue en $x_{0}$. C'est notamment le cas si $f\in C^{n}(I)$.

Fixons $x \in I$ avec $x\neq x_{0}$. Par Taylor-Lagrange a l'ordre $n-1$, il existe $c_{x}$ strictement entre $x$ et $x_{0}$ tel que
$$
f(x)=p_{n-1}(x)+\frac{f^{(n)}(c_{x})}{n!}(x-x_{0})^{n}.
$$
Par ailleurs,
$$
p_{n}(x)=p_{n-1}(x)+\frac{f^{(n)}(x_{0})}{n!}(x-x_{0})^{n},
$$
la difference donne
$$
f(x)-p_{n}(x)=[f^{(n)}(c_{x})-f^{(n)}(x_{0})] \frac{(x-x_{0})^{n}}{n!}.
$$
On a,
$$
\frac{f(x)-p_{n}(x)}{(x-x_{0})^{n}}=\frac{1}{n!}(f^{(n)}(c_{x})-f^{(n)}(x_{0})).
$$
Comme $c_{x}$ est entre $x$ et $x_{0}$, on a que $c_{x}\to x_{0}$ quand $x\to x_{0}$. Ainis, $f^{(n)}$ est continue et donc
$$
\lim_{ x \to x_{0} } \frac{f(x)-p_{n}(x)}{(x-x_{0})^{n}}=\frac{1}{n!}(f^{(n)}(\lim_{ x \to x_{0} } c_{x})-f^{(n)(x_{0})})=0 
$$
\end{proof}
